\chapter{File System Fundamentals}

\section{Introduction}

A file system is one of the most important components of an operating system. It provides a structured way to store, organize, retrieve, protect, and manage data on storage devices such as hard disks, SSDs, USB drives, and network storage systems.

For Software Development Engineer interviews, file systems are frequently discussed because they directly impact storage efficiency, application performance, reliability, security, and scalability.

\section{Learning Objectives}

After completing this chapter, you should be able to:

\begin{itemize}
\item Understand the purpose of file systems.
\item Explain file concepts and attributes.
\item Describe common file operations.
\item Differentiate file access methods.
\item Analyze directory structures.
\item Understand metadata and permissions.
\item Explain hard links and symbolic links.
\item Solve common interview questions on file systems.
\end{itemize}

%====================================================
\section{What is a File System}
%====================================================

\subsection{Definition}

A File System is a collection of methods and data structures that an operating system uses to store, organize, retrieve, and manage files on storage devices.

\subsection{Responsibilities}

\begin{itemize}
\item File Storage
\item File Retrieval
\item Directory Management
\item Access Control
\item Space Management
\item Metadata Maintenance
\item Data Integrity
\end{itemize}

\subsection{Examples}

\begin{table}[h]
\centering
\begin{tabular}{|l|l|}
\hline
Operating System & File System \\
\hline
Linux & ext4, XFS, Btrfs \\
Windows & NTFS, FAT32, exFAT \\
macOS & APFS, HFS+ \\
\hline
\end{tabular}
\caption{Common File Systems}
\end{table}

\section{File Concept}

\subsection{Definition}

A file is a named collection of related information stored on secondary storage.

\subsection{Examples}

\begin{itemize}
\item source.cpp
\item report.pdf
\item image.jpg
\item database.db
\item video.mp4
\end{itemize}

\subsection{Real-World Analogy}

Think of a file cabinet:

\begin{itemize}
\item Cabinet = Storage Device
\item Folder = Directory
\item Document = File
\end{itemize}

\section{File Attributes}

Every file contains metadata called file attributes.

\subsection{Common Attributes}

\begin{itemize}
\item Name
\item Identifier
\item Type
\item Location
\item Size
\item Creation Time
\item Modification Time
\item Access Time
\item Owner
\item Permissions
\end{itemize}

\subsection{Example}

Linux:

\begin{verbatim}
-rw-r--r-- 1 user user 2048 Jan 1 notes.txt
\end{verbatim}

Contains:

\begin{itemize}
\item Permissions
\item Owner
\item Size
\item Timestamp
\item Filename
\end{itemize}

%====================================================
\section{File Operations}
%====================================================

Common file operations include:

\begin{enumerate}
\item Create
\item Open
\item Read
\item Write
\item Append
\item Seek
\item Rename
\item Close
\item Delete
\item Truncate
\end{enumerate}

\subsection{Linux Example}

Create file:

\begin{verbatim}
touch notes.txt
\end{verbatim}

Delete file:

\begin{verbatim}
rm notes.txt
\end{verbatim}

\subsection{Windows Example}

Create file:

\begin{verbatim}
type nul > notes.txt
\end{verbatim}

Delete file:

\begin{verbatim}
del notes.txt
\end{verbatim}

%====================================================
\section{Access Methods}
%====================================================

File systems provide different ways to access data.

\begin{enumerate}
\item Sequential Access
\item Direct Access
\item Indexed Access
\end{enumerate}

%====================================================
\section{Sequential Access}
%====================================================

\subsection{Definition}

Records are accessed in sequence.

\subsection{Example}

Reading a text file line-by-line.

\begin{figure}[h]
\centering
\begin{tikzpicture}
\node[draw] (a) {Record 1};
\node[draw,right=1.5cm of a] (b) {Record 2};
\node[draw,right=1.5cm of b] (c) {Record 3};

\draw[->] (a)--(b);
\draw[->] (b)--(c);
\end{tikzpicture}
\caption{Sequential Access}
\end{figure}

\subsection{Advantages}

\begin{itemize}
\item Simple
\item Efficient for streaming
\end{itemize}

\subsection{Disadvantages}

\begin{itemize}
\item Slow random access
\end{itemize}

%====================================================
\section{Direct Access}
%====================================================

\subsection{Definition}

Any record can be accessed directly.

\subsection{Example}

Database record lookup.

\begin{figure}[h]
\centering
\begin{tikzpicture}
\node[draw] (cpu) {Request};
\node[draw,right=4cm of cpu] (rec) {Record 100};

\draw[->] (cpu)--(rec);
\end{tikzpicture}
\caption{Direct Access}
\end{figure}

\subsection{Advantages}

\begin{itemize}
\item Fast retrieval
\item Suitable for databases
\end{itemize}

\subsection{Disadvantages}

\begin{itemize}
\item Additional indexing overhead
\end{itemize}

%====================================================
\section{Indexed Access}
%====================================================

\subsection{Definition}

An index maps records to storage locations.

\subsection{Example}

Database B+ Tree index.

\begin{figure}[h]
\centering
\begin{tikzpicture}
\node[draw] (idx) {Index};

\node[draw,below left=1cm and 2cm of idx] (r1) {R1};
\node[draw,below=1cm of idx] (r2) {R2};
\node[draw,below right=1cm and 2cm of idx] (r3) {R3};

\draw[->] (idx)--(r1);
\draw[->] (idx)--(r2);
\draw[->] (idx)--(r3);
\end{tikzpicture}
\caption{Indexed Access}
\end{figure}

\subsection{Advantages}

\begin{itemize}
\item Efficient searching
\item Supports large datasets
\end{itemize}

\subsection{Disadvantages}

\begin{itemize}
\item Extra storage for index
\end{itemize}

%====================================================
\section{Directory Structures}
%====================================================

Directories organize files into a hierarchy.

Common structures:

\begin{enumerate}
\item Single-Level Directory
\item Two-Level Directory
\item Tree-Structured Directory
\item Acyclic Graph Directory
\item General Graph Directory
\end{enumerate}

%====================================================
\section{Single-Level Directory}
%====================================================

\subsection{Definition}

All files reside in one directory.

\begin{figure}[h]
\centering
\begin{tikzpicture}
\node[draw] (root) {/};

\node[draw,below left=1cm and 1cm of root] (f1) {A.txt};
\node[draw,below=1cm of root] (f2) {B.txt};
\node[draw,below right=1cm and 1cm of root] (f3) {C.txt};

\draw (root)--(f1);
\draw (root)--(f2);
\draw (root)--(f3);
\end{tikzpicture}
\caption{Single-Level Directory}
\end{figure}

\subsection{Advantages}

\begin{itemize}
\item Simple
\end{itemize}

\subsection{Disadvantages}

\begin{itemize}
\item Filename conflicts
\item Poor scalability
\end{itemize}

%====================================================
\section{Two-Level Directory}
%====================================================

\subsection{Definition}

Each user has a separate directory.

\begin{figure}[h]
\centering
\begin{tikzpicture}
\node[draw] (root) {Root};

\node[draw,below left=1cm and 2cm of root] (u1) {User1};
\node[draw,below right=1cm and 2cm of root] (u2) {User2};

\draw (root)--(u1);
\draw (root)--(u2);

\end{tikzpicture}
\caption{Two-Level Directory}
\end{figure}

\subsection{Advantages}

\begin{itemize}
\item Avoids name conflicts
\end{itemize}

\subsection{Disadvantages}

\begin{itemize}
\item Limited organization
\end{itemize}

%====================================================
\section{Tree-Structured Directory}
%====================================================

\subsection{Definition}

Directories may contain files and subdirectories.

\begin{figure}[h]
\centering
\begin{tikzpicture}

\node[draw] (root) {/};

\node[draw,below left=1cm and 2cm of root] (home) {home};
\node[draw,below right=1cm and 2cm of root] (etc) {etc};

\node[draw,below=1cm of home] (user) {user};

\node[draw,below=1cm of user] (docs) {documents};

\draw (root)--(home);
\draw (root)--(etc);
\draw (home)--(user);
\draw (user)--(docs);

\end{tikzpicture}
\caption{Tree-Structured Directory}
\end{figure}

\subsection{Advantages}

\begin{itemize}
\item Highly scalable
\item Natural hierarchy
\end{itemize}

\subsection{Disadvantages}

\begin{itemize}
\item More complex traversal
\end{itemize}

%====================================================
\section{Acyclic Graph Directories}
%====================================================

\subsection{Definition}

Allows shared files and directories without cycles.

\subsection{Advantages}

\begin{itemize}
\item Resource sharing
\item No infinite traversal loops
\end{itemize}

\subsection{Disadvantages}

\begin{itemize}
\item More complex management
\end{itemize}

%====================================================
\section{General Graph Directories}
%====================================================

\subsection{Definition}

Allows arbitrary links, including cycles.

\subsection{Advantages}

\begin{itemize}
\item Maximum flexibility
\end{itemize}

\subsection{Disadvantages}

\begin{itemize}
\item Infinite traversal risk
\item Garbage collection challenges
\end{itemize}

%====================================================
\section{File Metadata}
%====================================================

\subsection{Definition}

Metadata is information about a file rather than file content.

\subsection{Examples}

\begin{itemize}
\item Owner
\item Size
\item Permissions
\item Creation Time
\item Modification Time
\end{itemize}

Linux command:

\begin{verbatim}
stat file.txt
\end{verbatim}

%====================================================
\section{File Permissions}
%====================================================

\subsection{Linux Permissions}

\begin{verbatim}
-rwxr-xr--
\end{verbatim}

\subsection{Meaning}

\begin{table}[h]
\centering
\begin{tabular}{|c|c|}
\hline
Symbol & Meaning \\
\hline
r & Read \\
w & Write \\
x & Execute \\
\hline
\end{tabular}
\caption{Permission Symbols}
\end{table}

\subsection{Permission Groups}

\begin{itemize}
\item Owner
\item Group
\item Others
\end{itemize}

%====================================================
\section{Links}
%====================================================

A link is a directory entry referencing a file.

Types:

\begin{enumerate}
\item Hard Link
\item Symbolic Link
\end{enumerate}

%====================================================
\section{Hard Links}
%====================================================

\subsection{Definition}

A hard link is another name for the same inode.

\subsection{Linux Example}

\begin{verbatim}
ln file1 file2
\end{verbatim}

\subsection{Properties}

\begin{itemize}
\item Same inode
\item Cannot span file systems
\item Usually not allowed for directories
\end{itemize}

\subsection{Illustration}

\begin{figure}[h]
\centering
\begin{tikzpicture}

\node[draw] (f1) {file1};
\node[draw,right=4cm of f1] (inode) {inode 100};

\node[draw,below=1cm of f1] (f2) {file2};

\draw[->] (f1)--(inode);
\draw[->] (f2)--(inode);

\end{tikzpicture}
\caption{Hard Links}
\end{figure}

%====================================================
\section{Symbolic Links}
%====================================================

\subsection{Definition}

A symbolic link stores a pathname pointing to another file.

\subsection{Linux Example}

\begin{verbatim}
ln -s target link
\end{verbatim}

\subsection{Properties}

\begin{itemize}
\item Different inode
\item Can span file systems
\item Can point to directories
\item Can become dangling
\end{itemize}

\subsection{Illustration}

\begin{figure}[h]
\centering
\begin{tikzpicture}

\node[draw] (link) {soft link};

\node[draw,right=5cm of link] (target) {target file};

\draw[->] (link)--(target);

\end{tikzpicture}
\caption{Symbolic Link}
\end{figure}

%====================================================
\section{Hard Link vs Soft Link}
%====================================================

\begin{table}[h]
\centering
\small
\begin{tabular}{|p{4cm}|p{5cm}|p{5cm}|}
\hline
Feature & Hard Link & Soft Link \\
\hline
References & Inode & Pathname \\
\hline
Same Inode & Yes & No \\
\hline
Cross Filesystems & No & Yes \\
\hline
Directory Support & Usually No & Yes \\
\hline
Dangling Possible & No & Yes \\
\hline
Delete Original & Still Works & Breaks Link \\
\hline
\end{tabular}
\caption{Hard Link vs Symbolic Link}
\end{table}

%====================================================
\section{Linux Examples}
%====================================================

List files:

\begin{verbatim}
ls
\end{verbatim}

Detailed listing:

\begin{verbatim}
ls -l
\end{verbatim}

Create directory:

\begin{verbatim}
mkdir project
\end{verbatim}

Show metadata:

\begin{verbatim}
stat file.txt
\end{verbatim}

Create hard link:

\begin{verbatim}
ln a.txt b.txt
\end{verbatim}

Create symbolic link:

\begin{verbatim}
ln -s a.txt link.txt
\end{verbatim}

%====================================================
\section{Windows Examples}
%====================================================

Create directory:

\begin{verbatim}
mkdir project
\end{verbatim}

List files:

\begin{verbatim}
dir
\end{verbatim}

Create symbolic link:

\begin{verbatim}
mklink link.txt target.txt
\end{verbatim}

Create directory junction:

\begin{verbatim}
mklink /J linkdir targetdir
\end{verbatim}

%====================================================
\section{Directory Structure Comparison}
%====================================================

\begin{table}[h]
\centering
\small
\begin{tabular}{|p{4cm}|p{4cm}|p{5cm}|}
\hline
Structure & Advantage & Disadvantage \\
\hline
Single-Level & Simple & Name Conflicts \\
\hline
Two-Level & User Isolation & Limited Hierarchy \\
\hline
Tree & Scalable & More Complex \\
\hline
Acyclic Graph & Sharing & Management Overhead \\
\hline
General Graph & Flexible & Cycles Possible \\
\hline
\end{tabular}
\caption{Directory Structure Comparison}
\end{table}

%====================================================
\section{Interview Focus}
%====================================================

\subsection{Frequently Asked Topics}

\begin{itemize}
\item Hard Link vs Soft Link
\item Directory Structures
\item Inodes
\item Metadata
\item File Permissions
\item Access Methods
\end{itemize}

\subsection{Golden Rule}

\begin{center}
Hard Link = Same Inode

Soft Link = Path Reference
\end{center}

%====================================================
\section{Key Takeaways}
%====================================================

\begin{itemize}
\item File systems organize and manage storage.
\item Files are named collections of data.
\item Metadata describes file properties.
\item Sequential, direct, and indexed access are common methods.
\item Tree directories dominate modern systems.
\item Hard links share the same inode.
\item Symbolic links store paths.
\item Linux and Windows support hierarchical directories.
\item Permissions provide security.
\item Directory design affects scalability.
\end{itemize}

%====================================================
\section{20 Interview Questions with Answers}
%====================================================

\begin{enumerate}
\item What is a file system?\\
Answer: A mechanism for organizing and managing files on storage devices.

\item What is a file?\\
Answer: A named collection of related data.

\item What are file attributes?\\
Answer: Metadata such as size, owner, permissions, and timestamps.

\item What is metadata?\\
Answer: Information about a file.

\item What is sequential access?\\
Answer: Accessing records in order.

\item What is direct access?\\
Answer: Accessing any record directly.

\item What is indexed access?\\
Answer: Using an index for retrieval.

\item What is a directory?\\
Answer: A structure that organizes files.

\item What is a single-level directory?\\
Answer: All files stored in one directory.

\item What is a two-level directory?\\
Answer: Separate directories per user.

\item What is a tree directory?\\
Answer: Hierarchical directory structure.

\item What is an acyclic graph directory?\\
Answer: Shared files without cycles.

\item What is a general graph directory?\\
Answer: Directory structure allowing cycles.

\item What are file permissions?\\
Answer: Access control settings.

\item What does chmod modify?\\
Answer: File permissions.

\item What is a hard link?\\
Answer: Another name for the same inode.

\item What is a symbolic link?\\
Answer: A path reference to another file.

\item Can hard links cross file systems?\\
Answer: No.

\item Can symbolic links cross file systems?\\
Answer: Yes.

\item Which directory structure is most common?\\
Answer: Tree-structured directory.
\end{enumerate}

%====================================================
\section{20 MCQs with Answers}
%====================================================

\begin{enumerate}
\item A file system manages:
A) CPU
B) Memory
C) Storage
D) Cache

Answer: C

\item File metadata contains:
Answer: File Attributes

\item Sequential access reads:
Answer: In Order

\item Direct access allows:
Answer: Random Retrieval

\item Indexed access uses:
Answer: Index Structure

\item Tree directories are:
Answer: Hierarchical

\item Hard links share:
Answer: Inode

\item Symbolic links store:
Answer: Pathname

\item Linux file permissions include:
Answer: r, w, x

\item chmod changes:
Answer: Permissions

\item Hard links can cross filesystems:
Answer: No

\item Soft links can cross filesystems:
Answer: Yes

\item Dangling links occur with:
Answer: Symbolic Links

\item Common Linux filesystem:
Answer: ext4

\item Common Windows filesystem:
Answer: NTFS

\item File size is:
Answer: Metadata

\item Tree structure provides:
Answer: Scalability

\item Single-level directories suffer:
Answer: Name Conflicts

\item Acyclic graph allows:
Answer: Sharing

\item Most common directory structure:
Answer: Tree
\end{enumerate}

%====================================================
\section{Revision Notes}
%====================================================

\begin{itemize}
\item File systems organize storage.
\item Files contain data.
\item Metadata describes files.
\item Access methods include sequential, direct, and indexed.
\item Directories organize files.
\item Tree structure is most common.
\item Hard links share inodes.
\item Symbolic links store paths.
\item Permissions provide security.
\item Linux and Windows use hierarchical directories.
\end{itemize}

%====================================================
\section{Cheat Sheet}
%====================================================

\begin{center}
\begin{tabular}{|p{5cm}|p{8cm}|}
\hline
File & Named Collection of Data \\
\hline
Metadata & Data About Data \\
\hline
Sequential Access & Ordered Access \\
\hline
Direct Access & Random Access \\
\hline
Indexed Access & Index-Based Access \\
\hline
Single-Level Directory & One Directory \\
\hline
Two-Level Directory & Per User Directory \\
\hline
Tree Directory & Hierarchical Structure \\
\hline
Acyclic Graph & Sharing Without Cycles \\
\hline
General Graph & Cycles Allowed \\
\hline
Hard Link & Same Inode \\
\hline
Symbolic Link & Path Reference \\
\hline
Hard Link Cross FS & No \\
\hline
Soft Link Cross FS & Yes \\
\hline
Permissions & r, w, x \\
\hline
Linux FS & ext4 \\
\hline
Windows FS & NTFS \\
\hline
\end{tabular}
\end{center}